\chapter{ABI and Cross-Platform Compilation}

\section{Introduction}

We've seen how C code is compiled, linked, and executed. But what happens when you need to run your compiled binary on a different machine with a different CPU architecture or operating system? Why can't you just copy the binary and run it?

This chapter explores \textbf{ABI (Application Binary Interface)}—the set of conventions that enable binary code to work together at the machine level. Unlike API (source-level compatibility), ABI is about binary compatibility: how functions are called, how data is laid out in memory, and how the operating system expects programs to behave.

Understanding ABI is crucial for:
\begin{itemize}
    \item \textbf{Cross-platform development}: Writing code that compiles and runs correctly on different platforms
    \item \textbf{Cross-compilation}: Building binaries for one architecture on another
    \item \textbf{Library development}: Creating libraries that work across different compilers and platforms
    \item \textbf{Debugging}: Understanding mysterious crashes when mixing code from different sources
\end{itemize}

\section{What is an ABI?}

An \textbf{Application Binary Interface (ABI)} is a contract between compiled code components that specifies:

\begin{enumerate}
    \item \textbf{Calling convention}: How functions call each other (argument passing, return values, stack cleanup)
    \item \textbf{Data representation}: Size and alignment of data types
    \item \textbf{Register usage}: Which registers are preserved vs. volatile
    \item \textbf{Object file format}: ELF, PE, Mach-O, etc.
    \item \textbf{System call interface}: How to request kernel services
    \item \textbf{Runtime behavior}: Stack layout, exception handling, initialization
\end{enumerate}

\textbf{ABI vs. API}:

\begin{description}
    \item[API (Application Programming Interface)] Source-level interface. Defines function names, parameters, return types. Enables recompilation with different compilers/languages. Example: C header files.
    \item[ABI (Application Binary Interface)] Binary-level interface. Defines binary representation and calling conventions. Enables linking code compiled by different compilers. Example: System V AMD64 ABI.
\end{description}

\begin{tip}
API compatibility means you can recompile the code. ABI compatibility means you can link the binaries directly without recompilation.
\end{tip}

\section{Components of an ABI}

\subsection{Data Type Representation}

ABIs specify the size and representation of fundamental data types.

\textbf{Example: x86-64 Linux (LP64)}

\begin{tabular}{|l|c|c|l|}
\hline
\textbf{Type} & \textbf{Size} & \textbf{Alignment} & \textbf{Notes} \\
\hline
\texttt{char} & 1 byte & 1 & - \\
\hline
\texttt{short} & 2 bytes & 2 & - \\
\hline
\texttt{int} & 4 bytes & 4 & - \\
\hline
\texttt{long} & 8 bytes & 8 & Longs are 64-bit! \\
\hline
\texttt{long long} & 8 bytes & 8 & - \\
\hline
\texttt{void*} & 8 bytes & 8 & Pointers are 64-bit \\
\hline
\texttt{float} & 4 bytes & 4 & IEEE 754 single \\
\hline
\texttt{double} & 8 bytes & 8 & IEEE 754 double \\
\hline
\end{tabular}

\textbf{Naming schemes}:
\begin{itemize}
    \item \textbf{ILP32}: int, long, pointer are 32-bit
    \item \textbf{LP64}: long and pointer are 64-bit (int is 32-bit)
    \item \textbf{LLP64}: long long and pointer are 64-bit (Windows x64)
\end{itemize}

\subsection{Calling Conventions}

ABIs define how functions call each other—where arguments go, where return values go, who cleans the stack.

\textbf{System V AMD64 ABI (Linux/Unix x86-64)}:

Arguments (in order):
\begin{enumerate}
    \item RDI, RSI, RDX, RCX, R8, R9 (first 6 integer/pointer arguments)
    \item XMM0-XMM7 (first 8 floating-point arguments)
    \item Stack (remaining arguments)
\end{enumerate}

Return value:
\begin{itemize}
    \item Integer/pointer: RAX
    \item Floating-point: XMM0
    \item Large structs: Hidden pointer parameter
\end{itemize}

\textbf{Microsoft x64 ABI (Windows x64)}:

Arguments (in order):
\begin{enumerate}
    \item RCX, RDX, R8, R9 (first 4 arguments)
    \item Stack (remaining arguments)
\end{enumerate}

Return value:
\begin{itemize}
    \item RAX (integer/pointer)
    \item XMM0 (floating-point)
\end{itemize}

\textbf{Key difference}: Different registers for arguments! This is why you can't link object files between Windows and Linux.

\section{Cross-Compilation}

\subsection{What is Cross-Compilation?}

\textbf{Native compilation}: Build binaries for the same architecture/host
\begin{itemize}
    \item x86-64 Linux $\rightarrow$ x86-64 Linux binary
\end{itemize}

\textbf{Cross-compilation}: Build binaries for a different architecture/host
\begin{itemize}
    \item x86-64 Linux $\rightarrow$ ARM64 Linux binary
    \item x86-64 macOS $\rightarrow$ x86-64 Linux binary
\end{itemize}

\subsection{Cross-Compilation Toolchains}

A cross-compilation toolchain includes:
\begin{itemize}
    \item \textbf{Cross-compiler}: gcc that generates code for target architecture
    \item \textbf{Cross-assembler}: as for target
    \item \textbf{Cross-linker}: ld for target
    \item \textbf{Sysroot}: Target system's headers and libraries
\end{itemize}

\textbf{Naming convention}: \texttt{\{host\}-\{target\}-\{tool\}}

Examples:
\begin{itemize}
    \item \texttt{x86\_64-linux-gnu-gcc}: Native compiler for x86-64 Linux
    \item \texttt{aarch64-linux-gnu-gcc}: Cross-compiler for ARM64 Linux (from x86-64 host)
    \item \texttt{x86\_64-w64-mingw32-gcc}: Cross-compiler for Windows (from Linux host)
\end{itemize}

\subsection{Installing Cross-Compilers}

\textbf{Ubuntu/Debian}:
\begin{lstlisting}[language=bash]
# ARM64 cross-compiler
sudo apt-get install gcc-aarch64-linux-gnu

# Windows cross-compiler
sudo apt-get install gcc-mingw-w64-x86-64

# Check installation
aarch64-linux-gnu-gcc --version
x86_64-w64-mingw32-gcc --version
\end{lstlisting}

\subsection{Cross-Compiling a Program}

\textbf{Simple cross-compilation}:
\begin{lstlisting}[language=bash]
# Compile for ARM64 instead of x86-64
aarch64-linux-gnu-gcc -o hello_arm64 hello.c

# Check the output file
file hello_arm64
# Output: ELF 64-bit LSB executable, ARM aarch64...

# It won't run on x86-64!
./hello_arm64
# bash: ./hello_arm64: cannot execute binary file

# But it works on ARM64 systems
scp hello_arm64 user@arm64-machine:~
\end{lstlisting}

\section{ELF vs. PE vs. Mach-O}

Different operating systems use different executable file formats:

\subsection{ELF (Executable and Linkable Format)}

\textbf{Used by}: Linux, most Unix-like systems, many embedded systems

\textbf{File type}: \texttt{ELF 64-bit LSB executable, x86-64, version 1 (SYSV)}

\subsection{PE (Portable Executable)}

\textbf{Used by}: Windows

\textbf{File type}: \texttt{PE32+ executable (console) x86-64, for MS Windows}

\subsection{Mach-O}

\textbf{Used by}: macOS, iOS

\textbf{File type}: \texttt{Mach-O 64-bit executable x86\_64}

\section{Key Takeaways}

\begin{enumerate}
    \item \textbf{ABI is binary-level contract}: Unlike API (source-level), ABI defines how binaries interact—data layout, calling conventions, register usage.
    \item \textbf{Data types vary by platform}: \texttt{long}, \texttt{void*}, and even \texttt{int} may have different sizes. Use fixed-width types (\texttt{int32\_t}) for binary data.
    \item \textbf{Calling conventions differ}: x86-64 System V (Linux) uses different argument registers than Microsoft x64 (Windows), preventing binary compatibility.
    \item \textbf{Cross-compilation requires toolchains}: Build binaries for different architectures using cross-compilers like \texttt{aarch64-linux-gnu-gcc}.
    \item \textbf{Executable formats differ}: ELF (Linux), PE (Windows), and Mach-O (macOS) are incompatible at the binary level.
\end{enumerate}

\section{Looking Ahead}

In Chapter 10, we'll explore \textbf{process loading and memory layout}—how the operating system loads executables into memory and sets up the runtime environment.
